\subsection*{略語}

\addcontentsline{toc}{subsection}{略語}

\begin{center}
\small
\par\smallskip\noindent\refstepcounter{table}\label{tab:ch14-02}表 \thetable: 略語の整理\par\smallskip
表 \thetable は、略語の整理を比較・確認しやすい形で整理したものである。\par\smallskip
\begin{longtable}{p{0.17\linewidth}p{0.78\linewidth}}
\toprule
略語 & 日本語での意味 \\
\midrule
\endfirsthead
\toprule
略語 & 日本語での意味 \\
\midrule
\endhead
ACM & 計算機協会。情報処理分野の国際的な学会・専門職団体である。 \\
\midrule
CCPA & カリフォルニア州消費者プライバシー法。データプライバシー法の例である。 \\
\midrule
EEA & 欧州経済領域。 \\
\midrule
ENAEE & 欧州工学教育認定ネットワーク。工学教育認定に関わる組織である。 \\
\midrule
EU & 欧州連合。 \\
\midrule
GDPR & 一般データ保護規則。EU・EEAにおける個人データ保護規則の代表例である。 \\
\midrule
IEA & 国際エンジニアリング連合。工学教育・専門職認定に関わる国際的枠組みである。 \\
\midrule
IEC & 国際電気標準会議。電気・電子・関連技術の国際標準化組織である。 \\
\midrule
IEEE CS & IEEEコンピュータソサエティ。 \\
\midrule
IFIP & 国際情報処理連合。 \\
\midrule
IP & 知的財産。 \\
\midrule
ISO & 国際標準化機構。 \\
\midrule
NDA & 秘密保持契約。仕事を通じて知った機密情報を守る契約である。 \\
\midrule
UI/UX & 利用者インタフェース／利用者体験。画面や操作感だけでなく、利用者が受け取る体験を含む。 \\
\midrule
WIPO & 世界知的所有権機関。 \\
\midrule
WTO & 世界貿易機関。 \\
\midrule
\bottomrule
\end{longtable}
\normalsize
\end{center}
